\documentclass[aspectratio=169, 12pt]{beamer}

% --- Language and Fonts ---
\usepackage{fontspec}
\usepackage[english]{babel}
\usepackage{amsmath, amssymb, bm}
% tikz
\usepackage{tikz}
\usetikzlibrary{shapes, arrows.meta, positioning, shadows, calc}

% --- Bibliography (BibTeX + natbib) ---
\usepackage{natbib}
\bibliographystyle{chicago} 

% =========================================
% ===== keybox ====================
% =========================================
\usepackage[most]{tcolorbox}
\newtcolorbox{keybox}[1]{
  enhanced,
  colback=headbg!4,
  colframe=headbg,
  coltitle=headfg,
  colbacktitle=headbg,
  title=#1,
  fonttitle=\bfseries,
  boxrule=0.8pt,
  arc=1.5mm,
  left=2mm,
  right=2mm,
  top=1mm,
  bottom=1mm,
  before skip=0.2cm,
  after skip=0.2cm
}

\usepackage{booktabs}

% =========================================
% ===== Custom Footline (Page Number) =====
% =========================================


\setbeamercolor{footline}{fg=headbg}
\setbeamerfont{footline}{size=\scriptsize}

\setbeamertemplate{footline}{%
  \hfill%
  \usebeamercolor[fg]{footline}%
  \usebeamerfont{footline}%
  \insertframenumber\,/\,\inserttotalframenumber\kern1em\vskip3pt%
}

\setbeamertemplate{navigation symbols}{}

\usetheme{default}
\usepackage{xcolor}

\definecolor{headbg}{RGB}{30, 60, 120}
\definecolor{headfg}{RGB}{255, 255, 255}
\definecolor{accent}{RGB}{200, 50, 50}

\setbeamercolor{headline}{bg=headbg, fg=headfg}
\setbeamercolor{section in head/foot}{bg=headbg, fg=headfg}
\setbeamercolor{subsection in head/foot}{bg=headbg!80, fg=headfg}
\setbeamercolor{frametitle}{bg=headbg, fg=headfg}

% Step indicator command
\newcommand{\stepindicator}[1]{%
  \hfill {\scriptsize \textit{\textcolor{headbg}{Step #1 / 6}}}%
}
\DeclareMathOperator*{\argmin}{arg\,min}

\usepackage{hyperref}


% --- Title ---
\title[]{A First Course in Causal Inference Ch.~9}
\author[Kodai Minesaki]{Kodai Minesaki}
\institute[Faculty of Engineering, The University of Tokyo]{Faculty of Engineering, The University of Tokyo}
\date{\today}

\begin{document}

% --- Title Slide ---
\begin{frame}
\titlepage
\end{frame}


\begin{frame}

  {
    \large
    Ch.~9 Bridging Finite and Super Population Causal Inference
  }
  
\end{frame}

\begin{frame}{Ch.~9 Bridging Finite and Super Population Causal Inference}
  \begin{keybox}{Definition~9.1 internal validity}
    The statistical analysis is valid for the study samples at hand.
  \end{keybox}
  \begin{keybox}{Definition~9.2 external validity}
    The statistical analysis is valid for a broader population beyond the study samples. 
  \end{keybox}
  \begin{itemize}
    \item Do finite population results generalize to a bigger population?
    \item The external validity of the results depends on the sampling process of the units
    \begin{itemize}
      \item If the finite population represents a sample of a larger population, the results also have external validity.
    \end{itemize}
    \item The results based on randomization inference may not generalize; it ensures only internal validity. 
  \end{itemize}
  
\end{frame}

\begin{frame}{Ch.~9 Bridging Finite and Super Population}
  To rigorously discuss this issue, we introduce a frmework with two levels of randomness:
  \begin{itemize}
    \item Randomness due to the sampling process of the units;
    \item Randomness due to the random assignment of the treatment.
  \end{itemize}
  Actually, this framework does not really solve the problem mentioned above, but it is beneficial for the following two purposes:
  \begin{itemize}
    \item it gives a different perspective on randomized experiments;
    \item it serves as a bridge between experimental studies and observational studies.
  \end{itemize}
\end{frame}

\begin{frame}{9.1 CRE}
  
\end{frame}

\begin{frame}{9.1 CRE}
  
\end{frame}

\begin{frame}{9.1 CRE}
  
\end{frame}

\begin{frame}{9.1 CRE}
  
\end{frame}

\begin{frame}{9.1 CRE}
  
\end{frame}

\begin{frame}{9.1 CRE}
  
\end{frame}

\begin{frame}{9.1 CRE}
  
\end{frame}



% =========================

% =========================
\section{}
\begin{frame}[allowframebreaks, noframenumbering]{References}
  \beamerdefaultoverlayspecification{}
  \small
  \bibliography{../references}
\end{frame}

\end{document}